% !TeX root = ../main.tex
\section*{How to read this document}
\addcontentsline{toc}{section}{How to read this document}

\paragraph{Purpose.} Coding agents are building a tone estimator for Guitar
Virtual Amp (\file{guitar_core/src/tone/}). This document explains the problem,
the research it builds on, the design that was chosen, exactly what is
implemented, and how well it works, so the owner can follow and steer that
work. It replaces an earlier conference-style literature review and keeps all of
its references.

\paragraph{Structure.}
\begin{itemize}
\item \textbf{Part I, Background} (Sections~\ref{sec:problem}--\ref{sec:search}):
  what a tone is, the signal processing behind a guitar rig, how two sounds are
  compared, and how a search finds settings. Read this first if any of
  ``impulse response'', ``minimum phase'' or ``CMA-ES'' is unfamiliar.
\item \textbf{Part II, The literature} (Sections~\ref{sec:datasets}--\ref{sec:comparison}):
  datasets, effect recognition, amplifier modelling, working without the dry
  signal, style transfer and sound matching, commercial products, and a
  comparison with failure modes.
\item \textbf{Part III, Guitar Virtual Amp} (Sections~\ref{sec:design}--\ref{sec:next}):
  the design, the implementation component by component (with the formulas the
  code uses), the benchmark and its results, and the next steps.
\item \textbf{Appendices}: glossary; code map and commands.
\end{itemize}

\paragraph{Labels.} Coloured boxes: \emph{Key idea} (intuition), \emph{In Guitar
Virtual Amp} (what the code does), \emph{Caution} (pitfalls), \emph{Worked
example}, \emph{Next steps}. Inline tags: \measured{} run by us in this
repository; \reported{} stated by a paper's authors; \expected{} a hypothesis.
For the app: \done{} in the committed code, \wip{} in the working tree or being
changed, \planned{} only in the roadmap.

\paragraph{Snapshot and moving parts.} Written on 25 November 2026, updated on
26 November. The tone module was committed in \code{7db5015} (search, pluck
stand-in, match IR), \code{25f53ea} (measured delay, separated-stem targets, the
tone stage of the transcription pipeline) and \code{9939922} (preset prior, two
saved presets). The 26 November update adds a breakdown of where the error comes
from and a colour calibration of the stand-in, measured in the benchmark but not
yet used by the app (Section~\ref{sec:results-colour}). The results in
Section~\ref{sec:results} are from the committed defaults.

\paragraph{Related files.} \file{ROADMAP.md} (``Tone from the recording''),
\file{CLAUDE.md} (section ``Tone from a recording''), and the companion
document on tablature, \file{docs/papers/tablature-transcription/}, whose
separation and note transcription this work reuses~\cite{deangelo2026tabreview}.

\subsection*{The whole story in one table}
\addcontentsline{toc}{subsection}{The whole story in one table}

\begin{table}[h]
\centering
\caption{Summary. Distances are in dB of log-band difference (lower is better,
Section~\ref{sec:app-distance}).}
\label{tab:summary}
\small
\begin{tabularx}{\textwidth}{@{}>{\raggedright}p{2.6cm}>{\raggedright}p{4.4cm}X@{}}
\toprule
Question & Literature & Guitar Virtual Amp \\
\midrule
Which effects are present? & Close to solved on synthetic isolated guitar (set F1 0.987 on 3-effect chains) \reported & Not estimated separately: the search switches pedals off \\
Knob values of a known chain? & Regression MAE 0.02--0.05 on a unit scale for known plugins \reported; search by re-rendering beats prediction (wet SI-SDR 18.2 $\to$ 23.1\dB) \reported & CMA-ES search through the app's own renderer \done \\
Without the dry signal? & Unpaired amp models score 77--84 MUSHRA \reported; effect removal; no notes-based resynthesis found & Dry signal synthesised from the transcribed notes (plucked-string model) \done; its colour calibrated on the player's DI measured \wip \\
On a stem from a commercial mix? & Only one measurement: effect retrieval halves on separated stems \reported & Separated-stem targets measured: search and preset + match tie (5.9 and 5.8\dB) \measured \\
Final colour (cabinet, mix EQ) & Spectral matching to an IR, after the gain is right (Tone Match) & Minimum-phase match IR in a Cabinet pedal \done \\
How good is it? & No benchmark for the full task & EGDB, 12 items: 5.6--6.5\dB{} against 5.8--6.9\dB{} for preset + match; 1.7--3.1\dB{} with the true DI; 5.0--5.9\dB{} with a calibrated stand-in \measured \\
\bottomrule
\end{tabularx}
\end{table}
